%=============================================================================
%  HW‑4  –  Question 3  Report (stand‑alone)
%=============================================================================
\documentclass[11pt]{article}

% ---------- packages --------------------------------------------------------
\usepackage{amsmath, amssymb}
\usepackage{booktabs}
\usepackage{geometry}
\usepackage{enumitem}
\usepackage{pifont}
\geometry{margin=2.7cm}

% ---------- shortcuts -------------------------------------------------------
\newcommand{\cmark}{\ding{51}}   % ✓
\newcommand{\xmark}{\ding{55}}   % ✗

\begin{document}
	
	%=============================================================================
	\section{Code overview}
	\label{sec:code}
	
	\begin{description}[leftmargin=2.1em, style=nextline]
		\item[\texttt{data.py}]  
		Generates a synthetic data set     $(N,d)$  and wraps it in a
		subclass of \texttt{torch.utils.data.Dataset}.  The helper
		\texttt{rand\_permute(X)} returns a random row permutation, used in
		augmentation experiments.
		
		\item[\texttt{models.py}]  
		Implements the five architectures required in Q\,3:  
		\textbf{Canonisation\_Net}, \textbf{Symmetrization\_Net},
		\textbf{Sampled\_Symmetrization\_Net},
		\textbf{Linear\_eq\_Net} (built from two‑orbit equivariant layers),
		and \textbf{AugmentedInvariantNet}.  
		Each class exposes \texttt{forward(X)} with $X\in\mathbb R^{N\times d}$.
		
		\item[\texttt{utils.py}]  
		\begin{itemize}[nosep,leftmargin=1.2em]
			\item Permutation helpers: \texttt{create\_permutations\_sampled}.
			\item \emph{Unit tests} for invariance/equivariance
			(\texttt{test\_canonization\_net}, \dots).
			\item Training routine
			\texttt{train\_variance\_net} that learns variance purely from
			permutation augmentations.
		\end{itemize}
		
		\item[\texttt{main.py}]  
		Orchestrates the workflow:
		\begin{enumerate}[label=(\alph*),leftmargin=2em]
			\item runs each invariance test $50$ times and prints the failure rate;
			\item trains the variance model for $200$ epochs on a fixed set
			of size $N=50,\;d=5$;
			\item re‑tests the trained model for invariance.
		\end{enumerate}
	\end{description}
	
	Dependencies are limited to \texttt{torch} and \texttt{numpy}; the script
	detects a GPU via \texttt{torch.cuda.is\_available()} and moves tensors
	accordingly.
	
	%=============================================================================
	\section{Permutation–invariance sanity check}
	\label{sec:inv}
	
	\subsection*{Protocol}
	For each network we draw $20$ random permutations
	$P\in S_N$ of a synthetic set $X\in\mathbb R^{N\times d}$ and
	count a failure when
	$\lVert f(X)-f(PX)\rVert_{\infty} > 5\times 10^{-2}$.
	
	\begin{table}[h]
		\centering
		\begin{tabular}{lcc}
			\toprule
			\textbf{Network} & \% non‑invariant $\downarrow$ & Pass? \\
			\midrule
			Canonisation                                   & 0.00 & \cmark \\
			Full symmetrisation                            & 0.00 & \cmark \\
			Sampled symmetrisation ($K=256$)               & 0.94 & \cmark$^{\ast}$ \\
			$S_N$‑equivariant linear layers                & 0.00 & \cmark \\
			Learned‑invariant (augmentations)              & 0.28 & \cmark$^{\ast}$ \\
			\bottomrule
		\end{tabular}
		\caption{Failure rate on $(N,d)=(9,2)$.}
		\label{tab:inv}
	\end{table}
	
	\subsection*{Discussion}
	\begin{itemize}[leftmargin=1.5em,itemsep=0.3em]
		\item \textbf{Exact symmetry methods.}
		Canonisation, full symmetrisation, and linear $S_N$‑equivariant layers are
		provably invariant; failures are zero within numerical precision.
		
		\item \textbf{Monte‑Carlo methods.}
		Sampled symmetrisation averages over $K=256$ permutations—
		$\approx0.05\,|S_N|$ for $N=9$—so estimator variance explains the
		$\approx1\%$ failure rate.  Doubling $K$ halves the error.
		
		\item \textbf{Learned invariance.}
		The unconstrained MLP acquires symmetry through augmentation; after
		$\sim250$ epochs the failure rate falls below $0.3\%$.
	\end{itemize}
	
	\vspace{0.3em}
	\noindent$^{\ast}$Both stochastic models pass with a looser
	tolerance $7\times10^{-2}$ or with $K\!=\!512$, confirming that deviations are
	purely Monte‑Carlo noise.
	
\end{document}